\section{Séptima iteración de la Fase 5: Mejoras en la pantalla}
\label{sec:F5I7}

Antes del proceso de lectura del GWC, que implica añadir más botones en
la pantalla, he decidido mejorar primero la presentación de esta, haciéndola
más operativa. Mi idea es que el manejo sea lo más natural posible, y
ahora mismo están juntas las pantallas del simulador y del editor
gráfico, lo que no es demasiado intuitivo, y presento demasiada información
que no es necesario tener allí.

\subsection{Redefinir la pantalla}
\label{sec:F5I7Pantalla}

La idea es disponer de tres pestañas en la pantalla. La primera es la
del simulador: si se quiere jugar no hace falta nada más. La segunda es
la del editor gráfico: al eliminar de ella la pantalla del simulador
gano espacio, por lo que ahora el editor ocupa la parte izquierda de la
pantalla, y al pulsar en un segmento sus propiedades aparecen a la
derecha, en un sistema similar al de los editores gráficos. La tercera
es la lista de segmentos, cada uno con la imagen que tiene asociada.

\begin{itemize}
	\item \textbf{Simulador}: 
	Solo la pantalla del juego dentro de la carcasa de la máquina, con
	sus botones, y debajo la información de la máquina y del cartucho.
	
	\item \textbf{Editor}:
	El mapa de segmentos a la izquierda y, al pulsar en un segmento, su
	panel de propiedades a la derecha. Debajo, la misma información que
	en el simulador, ampliada con los datos del diseño.
	
	\item \textbf{Segmentos y Sprites}:
	La lista de todos los segmentos del cartucho y, junto a cada uno, la
	imagen asociada si la tiene. Sustituye a la antigua tira de sprites,
	que mostraba las imágenes y no los segmentos. Aparecerán todos los segmentos y sus imágenes asociadas, separados en sublistas:
	\begin{itemize}
		\item Segmentos sin imagen
		\item Segmentos visibles con imagen
		\item Segmentos no visibles con imagen (los que se han unido con otros)
		\item Imágenes sin segmento (si las hay)
	\end{itemize}
\end{itemize}

Las tres pestañas se definen en \texttt{index.html}. Es un cambio
solo de presentación: el motor, el puente y el módulo de cartuchos no
cambian.

\subsection{La barra de botones}
\label{sec:F5I7Botones}

La barra superior es única y común a las tres pestañas, pero cada
pestaña solo muestra los botones que tienen sentido en ella. Los que no
corresponden se ocultan, no se deshabilitan, para que la barra quede
limpia:

\begin{table}[H]
	\centering
	\begin{tabular}{lccc}
		\hline
		\textbf{Botón} & \textbf{Simulador} & \textbf{Editor} & \textbf{Segmentos y Sprites} \\
		\hline
		Cargar Setup        & Sí & Sí & No \\
		Cargar Cartucho\footnotemark & No & Sí & No \\
		Cargar GWI (ZIP)    & Sí & Sí & No \\
		Cargar GWC          & Sí & No & No \\
		Verificar y Guardar & No & Sí & No \\
		Generar GWI+GWC     & No & Sí & No \\
		Deshacer / Rehacer  & No & Sí & No \\
		\hline
	\end{tabular}
	\caption{Botones visibles en cada pestaña}
	\label{tab:F5I7Botones}
\end{table}
\footnotetext{Se cambió el nombre a \textbf{Cargar GWB} en la novena iteración, ver sección \ref{sec:F5I9Entrada}.}


\texttt{BaseDirector} decide qué botones se ven al cambiar de pestaña
y \texttt{estilos.css} los oculta. La línea de estado, donde aparecen
los avisos y los errores, queda siempre visible, sea cual sea la
pestaña.

\subsection{La información de cada pestaña}
\label{sec:F5I7Informacion}

El panel de información deja de ser común y \texttt{BaseDirector} pasa
a mostrar en cada pestaña solo lo que interesa en ella:

\begin{itemize}
	\item \textbf{Simulador}: el nombre del cartucho, la resolución de
	la pantalla y los botones que maneja la máquina (2, 4 o 5).
	
	\item \textbf{Editor}: lo mismo que en el simulador, añadiendo el
	número de segmentos, el número de imágenes del cartucho, el tamaño
	del GWI (aproximado) y el tamaño del GWC (si se han generado o cargado).

	\item \textbf{Segmentos y Sprites}: sin panel de información; la propia lista
	es la información.
\end{itemize}

El objetivo de los dos tamaños es ver lo que se gana al pasar del GWI
al GWC, no la ocupación en disco, por eso el del GWI no es el del ZIP
comprimido sino la suma de lo que contiene: el JSON más las imágenes
PNG, cada fichero con sus bytes reales y no con lo que ocupa en el
disco (un sprite de 978 bytes suma 978, aunque en el disco ocupe un
bloque de 1 KB). Es una aproximación: no hace falta que sea exacta, y
lo que añadan las cabeceras del PNG al guardarlo no se descuenta. El
del GWC sí es exacto, es el tamaño del fichero binario.

\subsection{La lista de segmentos}
\label{sec:F5I7Lista}

La lista la construye \texttt{BaseDescomponer}, sustituyendo a la tira
de sprites. Cada fila muestra el identificador del segmento, su tipo y
la imagen asociada. Cada segmento aparece en una de estas sublistas:

\begin{itemize}
	\item \textbf{Segmentos sin imagen:} los dígitos y los contadores,
	cuyo gráfico son las barras que los componen, que aparecen en la
	lista como segmentos propios, y el segmento de cartucho, que es la
	ficha de configuración y nunca se presenta.

	\item \textbf{Segmentos visibles con imagen:} los segmentos normales
	con su imagen, y las piezas de unión con la imagen combinada de sus
	miembros, que es la que se presenta en pantalla.

	\item \textbf{Segmentos no visibles con imagen:} los miembros que
	una unión ha absorbido, cada uno con su imagen original.

	\item \textbf{Imágenes sin segmento:} las que no utiliza ningún
	segmento, información que antes daba la tira de sprites y que así
	no se pierde. Si no hay ninguna, el apartado no aparece.
\end{itemize}

Al pulsar un segmento de la lista se pasa a la pestaña Editor con ese segmento seleccionado.

\subsection{Pestañas disponibles y partida en marcha}
\label{sec:F5I7Estados}

Las pestañas solo se pueden usar cuando tienen algo que mostrar, y lo
controla \texttt{BaseDirector}:

\begin{itemize}
	\item Sin setup cargado solo está disponible el Simulador.
	
	\item La pestaña «Editor» se habilita al cargar el setup, ya que es
	desde ella donde se carga la carpeta del cartucho base.

	\item La pestaña «Segmentos y Sprites» se habilita al cargar un
	cartucho editable, ya sea el base en su carpeta o el GWI en su ZIP.
	La carga del GWC se define en la siguiente iteración.
	
	\item Al pulsar Start la partida se juega en el Simulador. Mientras
	dura, las otras pestañas y los botones de carga quedan
	deshabilitados (no ocultos), y vuelven a estar disponibles al
	terminar la partida. Hasta ahora, al pulsar Start se ocultaban y no volvían a
	aparecer, y había que recargar la página.
\end{itemize}


\subsection{Pruebas}
\label{sec:F5I7Pruebas}

\begin{itemize}
	\item \textbf{Botones:} en cada pestaña se ven exactamente los de la
	Tabla~\ref{tab:F5I7Botones}, sin setup, con cartucho y con GWC.
	
	\item \textbf{Información:} el Simulador muestra nombre, resolución
	y botones, y el Editor añade segmentos, imágenes y los tamaños del
	GWI y del GWC, con los valores correctos del Parachute; el tamaño
	del GWI es la suma del JSON y de los PNG, no el del ZIP.

	\item \textbf{Lista:} cada segmento aparece en su sublista con su
	imagen, las piezas de unión con la imagen combinada, las imágenes
	sin segmento aparecen al final solo si las hay, y al pulsar un
	segmento se abre en el Editor.
	
	\item \textbf{Partida:} durante la partida no se puede cambiar de
	pestaña ni cargar; al terminar, todo vuelve a estar disponible sin
	recargar la página.
\end{itemize}